% Transcription — fonds Alexandre Grothendieck, shelfmark 56, batch 2 (pages 21-38).
% Established from the facsimile archives/batches/56/batch-02.pdf (a mirror of
% https://grothendieck.umontpellier.fr/56.pdf; no cover sheet, PDF page k =
% folder page 20+k), per the skill transcribe-grothendieck, with the 700 dpi
% tiles of archives/tiles/56/ and 400-900 dpi re-crops (turned 90° for page 27,
% 180° for the pencil at the foot of page 21, about 22° for the slanted note
% of page 23). The folder is thirty-eight pages; this is the second and last
% batch.
%
% Pages skipped: 33, a blank cream sheet; 38, the back cover of the offprint,
% which carries only its printed colophon (« Extrait des Comptes rendus des
% séances de l'Académie des Sciences, t. 255, p. 1366-1368, séance du
% 10 septembre 1962 », Gauthier-Villars) and the show-through of page 37 —
% the reference is given in the \note of page 35 instead.
%
% Summarised pages. (1) Pages 22, 24, 26, 28, 30 and 32 are, as in batch 1,
% leaves of his own typescript of « XII Tores maximaux, groupe de Weyl,
% sous-groupes de Cartan, centre réductif… » (SGA 3, Exp. XII), corrected by
% him in ink, on whose backs the manuscript notes are written; 22 and 28 were
% scanned upside down. Following batch 1 and the brief (his text, but
% published), each gets one French \note on what it states and his ink
% interventions are given in full; the typed text is not re-set. Page 30
% carries almost nothing of his (one inked bracket) and is kept only so that
% the run of typed leaves is not broken; this is a judgement call. The typed
% page numbers (23, 11 (?), 21, 12, 26, 31) show the leaves are not in the
% typescript's order. (2) Pages 35-37 are an offprint of someone else's
% printed Note aux Comptes rendus (1962) on Stiefel-Whitney classes of
% quadratic forms, the « tiré à part » of the inventory: one \note per page
% on what it states. No mark of his hand was found on it — the specks near
% « a.b » on page 37 and at the foot of page 35 are not writing — so nothing
% is transcribed from it. The author's name, printed in the heading of the
% Note, is a published byline and is given as such.
%
% What the batch is about. The manuscript notes continue batch 1's
% « Cohomologie des quadriques » (page 19, which broke off on the case n
% even of the Gysin sequence H^{n-2}(Q^{n-1})(-1) → H^n(Q^n) → H^n(U^n) → …,
% U^n = Q^n - Q^{n-1}).
%   21     the case n odd (H^{n-1}(Q^{n-1})(-1) an extension of H^{n+1}(Q^n)
%          by H^n(U^n)), then a programme a)-e): cohomology of Q^n zero in odd
%          degree; H^i(U^n) zero for i ≠ 0, n; dual statement; the maps between
%          H^i of Q^n and of Q^{n±1}, bijective / injective / surjective by
%          parity of n; H^{2i}(Q^n) of rank 1 except rank 2 for n even, i = n.
%   23     the cohomology ring: ξ^i ∈ H^{2i}(Q^n, μ(i)), λ_{2ν} ∈ H^{2ν}(Q^{2ν}),
%          the chain Q^0 ⊂ Q^1 ⊂ … ⊂ Q^4 with the bases in each degree, and the
%          boxed bases: Q^{2ν}: ξ^0 … ξ^{ν-1}, [ξ^ν, λ_ν], ½ξ^{ν+1} … ½ξ^{2ν};
%          Q^{2ν+1}: ξ^0 … ξ^ν, ½ξ^{ν+1} … ½ξ^{2ν+1}; with φ_ν = ½(ξ^ν + λ_ν).
%   25     the same with ξ_i = ½ξ^i: relations 2ξ_i = ξ^i, ξ^iξ_j = ξ_{i+j},
%          λ_ν^2 = 2ξ_{2ν} = ξ^{2ν}, and the functoriality for Q^i ⊂ Q^j:
%          restriction (λ_ν ↦ 0, φ_ν ↦ ξ_ν: « classe effective ») and Gysin
%          (λ_ν ⇝ 0, φ_ν ⇝ ξ_{ν+(j-i)}: « classe évanescente »).
%   27     (written sideways) the table of bases for Q^0 … Q^6.
%   29     the specialisation square for n = 2ν between H^n(Q^n), H^n(U^n),
%          H^{n-1}(U^{n-1})(1) and compact supports, with the images of
%          λ_ν, ξ^ν, φ_ν and generators ℓ_{2ν}, ℓ_{2ν+1}; N.B.: φ_ν defined by
%          a maximal isotropic subspace, α_ν + β_ν = ξ^ν, α_ν - β_ν = λ_ν,
%          φ_ν = α_ν. Two pencilled lines on derived tensor products at the top.
%   31     cup-product with ξ, H^{2i}(Q^{2ν}) → H^{2i+2}(Q^{2ν})(1), bijective
%          for i ≠ ν-1, ν (injective with free cokernel on φ_ν, resp. surjective
%          with free kernel on λ_ν); for Q^{2ν+1} bijective for i ≠ ν,
%          injective with cokernel Z/2Z for i = ν.
%   34     (tan paper, a fair hand) « Quadriques »: bases for linear
%          equivalence on Q^{2n} (1, H, …, H^{n-1}, S', S'', HS' = HS'', …,
%          H^n = S' + S'') and on Q^{2n+1} (1, H, …, H^n, H^{n+1}/2, …,
%          H^{2n+1}/2), with S', S'' the two families of linear subspaces.
%   35-37  the offprint (summarised).
%
% The hand of pages 21-31 is fast drafting in blue ink; the formulas carry
% and most connective words do not, and are left \ill{}. Page 34 reads
% better but its short words are still reduced to strokes.
%
% Notation, fixed as in batch 1: Q^n the smooth quadric of dimension n, U^n
% its affine part, ξ the hyperplane class, λ_ν the extra class in the middle
% degree, μ the roots of unity; ξ_i is his « ½ξ^i » (page 25), φ_ν his
% ½(ξ^ν + λ_ν), ℓ_i the generators he writes « l » on page 29.
%
% Readings to check first: the numbering « 2) » opening page 21 and the
% exponents of its first two lines; the domains of the maps in c) and c') on
% page 21; the slanted note above the box of page 23; « effective » and
% « évanescente » on page 25 and which case each belongs to; the labels on
% the arrows of page 29 and its N.B. (« directement », « isotrope maximal »);
% the word under S', S'' on page 34.
%
% Pass: Opus 5.5 (claude-opus-5-5), 2026-09-24 — first pass, unchecked against the pages by a human.

\documentclass[11pt,a4paper]{article}
% Le préambule commun à toutes les transcriptions du fonds Grothendieck.
%
% Il tient en une idée : **l'incertitude se compose, elle ne se résout pas.**
% Une transcription qui lisse un mot douteux a détruit la seule chose pour
% laquelle on la faisait — un lecteur doit pouvoir savoir, à chaque endroit,
% ce qui était sur le feuillet et ce qui a été ajouté. Les six macros qui
% suivent sont donc l'appareil critique, et non de la décoration.
%
% Les mêmes commandes sont reconnues par `scripts/render.mjs`, qui produit la
% vue de lecture du site. Une macro ajoutée ici doit l'être aussi là-bas,
% faute de quoi le rendu échouera bruyamment — ce qui est voulu : un rendu
% silencieusement faux est pire qu'un rendu absent.

\ProvidesPackage{grothendieck}[2026/08/08 Transcriptions du fonds Grothendieck]

\usepackage{amsmath,amssymb,amsthm}
\usepackage{mathrsfs}
\usepackage{stmaryrd}
% Les diagrammes commutatifs sont partout dans ce fonds : c'est la langue
% même de la géométrie algébrique de Grothendieck, et une transcription qui
% les rendrait en prose ne transcrirait rien.
\usepackage{tikz-cd}
% Français par défaut — c'est la langue des feuillets — mais l'anglais est
% chargé aussi : la traduction et le résumé sont des documents anglais, et la
% typographie française (espaces avant les deux-points) y serait une faute.
% Ces documents basculent par \selectlanguage{english} en tête de corps.
\usepackage[english,french]{babel}
\usepackage{fontspec}
\usepackage{geometry}
\geometry{a4paper,margin=25mm,marginparwidth=28mm}
\usepackage{marginnote}
\usepackage{xcolor}
% ulem, not soul. `soul`'s \st and \ul are built on a character-by-character
% scan that XeLaTeX's font machinery breaks: the document fails with an
% undefined control sequence rather than degrading. ulem does the same two jobs
% and is Unicode-engine safe. `normalem` stops it hijacking \emph, which the
% transcriptions use for Grothendieck's own emphasis.
\usepackage[normalem]{ulem}
\usepackage{hyperref}
\hypersetup{colorlinks=true,linkcolor=black,urlcolor=black,pdfborder={0 0 0}}

% --- Métadonnées du lot -----------------------------------------------------
% Reprises telles quelles par le script de rendu, qui les affiche en tête de
% la vue de lecture. Elles fixent ce que le fichier prétend couvrir : sans
% elles, une transcription flotte sans point d'attache dans un fonds de seize
% mille feuillets.

\newcommand{\folder}[1]{\gdef\grothFolder{#1}}
\newcommand{\batch}[1]{\gdef\grothBatch{#1}}
\newcommand{\leaves}[2]{\gdef\grothFirst{#1}\gdef\grothLast{#2}}
\newcommand{\foldertitle}[1]{\gdef\grothFoldertitle{#1}}
\gdef\grothFolder{?}\gdef\grothBatch{?}\gdef\grothFirst{?}\gdef\grothLast{?}\gdef\grothFoldertitle{}

% --- Le résumé --------------------------------------------------------------
%
% Il ouvre la lecture modernisée, sous le titre « Résumé », et il est composé
% en petit. Ce n'est pas un ornement : c'est le seul passage du document qui
% ne soit pas des mathématiques, et un lecteur doit pouvoir le reconnaître
% avant de l'avoir lu — puis le sauter s'il connaît déjà le sujet, ou s'y
% arrêter s'il ne le connaît pas. Le filet qui le suit marque la reprise du
% corps.

\newenvironment{resume}
  {\small\setlength{\parskip}{0.45em}}
  {\par\medskip\hrule\medskip}

% --- Mots-clés ---------------------------------------------------------------
%
% En anglais, en clôture du résumé de la lecture modernisée : le vocabulaire
% moderne sous lequel chercher ce que les feuillets construisent. Le manifeste
% du site (`npm run manifest`) les extrait pour étiqueter la cote entière —
% cette ligne est la source unique des tags, il n'y a pas de fichier à côté.
\newcommand{\keywords}[1]{\par\smallskip\noindent
  {\footnotesize\textsc{Keywords} --- \textit{#1}}\par}

% --- L'appareil critique ----------------------------------------------------

% Le numéro de feuillet, en marge. C'est l'ancre qui permet de régler un
% désaccord sur une formule en regardant *un* feuillet plutôt qu'une chemise
% de six cents. Le site s'en sert aussi pour tourner les pages du fac-similé
% quand on fait défiler la transcription.
\newcommand{\leaf}[1]{%
  \marginnote{\footnotesize\color{gray!70}\textsf{#1}}%
  \label{leaf:#1}\ignorespaces}

% Illisible : on ne devine pas. Un mot inventé qui se lit comme les autres est
% le pire résultat possible.
\newcommand{\ill}{\textcolor{red!60!black}{[\,\ldots\,]}}

% Lecture douteuse : le mot est donné, et signalé comme incertain.
\newcommand{\uncertain}[1]{\uline{#1}}

% Ajout de l'éditeur — une lettre manquante, un mot sous-entendu.
\newcommand{\add}[1]{\textcolor{blue!50!black}{[#1]}}

% Biffé par Grothendieck lui-même. Ce qu'il a rayé fait partie du document :
% une rature dit souvent où la pensée a changé de direction.
\newcommand{\struck}[1]{\sout{#1}}

% Note du transcripteur, sur l'état matériel du feuillet.
\newcommand{\note}[1]{\par\smallskip{\small\color{gray!60!black}\textit{[#1]}}\par\smallskip}

% Note marginale de Grothendieck, distincte du corps du texte.
\newcommand{\marginal}[1]{\par\smallskip\noindent\hspace{1em}%
  {\small\color{gray!60!black}#1}\par\smallskip}

% --- Titre ------------------------------------------------------------------

\AtBeginDocument{%
  \begin{center}
    {\large\bfseries Fonds Alexandre Grothendieck --- cote n\textsuperscript{o}~\grothFolder}\\[2pt]
    {\itshape\grothFoldertitle}\\[4pt]
    {\small feuillets \grothFirst--\grothLast\ (lot \grothBatch)}\\[2pt]
    {\footnotesize\color{gray!60!black}Transcription établie d'après le fac-similé numérisé
     par l'Université de Montpellier.\\
     Les lectures incertaines sont soulignées, les illisibles notées [\,\ldots\,],
     les ajouts entre crochets.}
  \end{center}
  \vspace{1em}}

\endinput


\folder{56}
\batch{2}
\pages{21}{38}
\dating{1962}
\watermark{Édition de démonstration}
\foldertitle{Lefschetz-Hodge - Vanishing cycles : notes manuscrites (s.d.), tiré à part (1962)}

\begin{document}

\section*{Cohomologie des quadriques (suite)}

\note{suite de la page 19 (lot 1), qui s'arrête sur le cas « 1) $n$ pair » ;
titre de l'édition}

\page{21}

\uncertain{2)} \emph{$n$ impair} $\geqslant 1$, la suite exacte \ill{} \ill{}
\ill{} :
\[
  0 \to H^{n}(U^{n}) \to H^{n-1}(Q^{n-1})(-1) \to H^{n+1}(Q^{n}) \to 0
\]
i.e. $H^{n-1}(Q^{n-1})$ est une extension de $H^{n+1}(Q^{n})$ par
\ill{}\note{les exposants de cette ligne sont écrits vite ; « $H^{n-1}(Q^{n-1})$ »
et « $H^{n+1}(Q^n)$ » sont lus d'après la suite exacte qui précède}

a) Cohomologie de $Q^n$ nulle en dim impaire.

b) $H^i(U^n)$ nul si $i \neq 0, n$, isomorphe à $A$ (resp. $A \otimes \nu_n$) si
$i = 0$ (resp. $i = n$)\note{« $i = 0$ (resp. $i = n$) » est entouré et relié
par un trait au « si » de la ligne}

b') Énoncé dual

c) $H^i(Q^{n+1}) \to H^i(Q^n)$ est
\[
  \left|
  \begin{array}{l}
    \text{bijectif si}\ \struck{\ill{}}\ i \geqslant 2,\ i \neq n \\
    \text{injectif pour tt}\ i
  \end{array}
  \right\}\ (n\ \text{pair})
  \qquad
  \left|
  \begin{array}{l}
    \text{bijectif si}\ i \geqslant 2,\ i \neq n+1 \\
    \text{surjectif pour tt}\ i \geqslant 1
  \end{array}
  \right\}\ (n\ \text{impair})
\]
\note{l'exposant du premier $Q$ est surchargé, et un « $(A)$ » biffé le suit ;
« $n+1$ » est une lecture douteuse. Dans la première accolade, le « $i \geqslant 2$ »
est écrit par-dessus « $2 \leqslant$ » et un mot biffé}

c') Énoncé dual : $H^i(Q^n) \to H^i(Q^{n+1})$ est
\[
  \left|
  \begin{array}{l}
    \text{bijectif si}\ i \leqslant 2n-2,\ i \neq n \\
    \text{injectif pour tt}\ i
  \end{array}
  \right\}\ n\ \text{pair}
  \qquad
  \left|
  \begin{array}{l}
    \text{bijectif si}\ i \leqslant 2n-2,\ i \neq n-1 \\
    \text{injectif pour tt}\ i \leqslant 2n-\ldots
  \end{array}
  \right\}\ n\ \text{impair}
\]
\note{en regard, à droite, un mot entouré, \ill{}, relié à la seconde
accolade, et au-dessous « \ill{} \ill{} » souligné ; sous la seconde ligne,
« \struck{\ill{}} \ill{} \ill{} » en interligne. L'exposant du second $Q$ est lu
$n+1$, douteux}

d) $H^{2i}(Q^n)$ \struck{est de rang} de rang $1$ \ill{} \ill{} pour $n$ pair et
\struck{\ill{}} $i = n$, où il est de rang $2$.

e) \ill{} Si $n$ pair, \ill{}

\marginal{$X_1 \subset X' = Y[t_1, \ldots, t_n]$, les deux s'envoyant sur $Y$}
\note{au crayon, tête-bêche au pied de la page, d'un trait plus fin ; rien ne
le relie à ce qui précède}

\page{22}

\note{feuillet de tapuscrit du même exposé que les pages impaires du lot 1,
paginé « 23 » : fin de la preuve que $G/Z$ n'a pas de sous-groupe de type
multiplicatif autre que le groupe unité ($H \cap G/Z_0 = 0$, donc
$H \to G/G_0$ injectif, $H$ radiciel et de rang majoré par celui de $G/G_0$,
d'où un plus grand sous-groupe de type multiplicatif $Z$) ; par Exp VI et
Exp VII, une extension commutative de deux groupes de type multiplicatif est de
type multiplicatif, d'où $Z'/Z = 0$ ; par le « principe de l'extension finie »,
il en est de même de $(G/Z)_K = G_K/Z_K$ ; puis le début de la preuve que $Z$ est
un centre réductif : tout $u : H \to G_S$, $H$ de type multiplicatif et de type
fini sur $S$, se factorise par $Z_S$, par réduction au cas d'un corps
(Exp IX 5.2 et 6.8)}

Sa main : les flèches et le « $\cap$ » ajoutés à l'encre ; « Cela
\struck{\ill{}} \add{implique} que la famille des sous-groupes de type
multiplicatifs de $G$ admet un plus grand élément », le mot étant écrit à
l'encre sur un mot tapé.

\page{23}

$\mathbf{P}^1 \times \mathbf{P}^1 \simeq Q^2 \subset Q^1 \simeq \mathbf{P}^1$
\qquad $\dfrac{n-1}{2} + 1 = \dfrac{n+1}{2}$\note{sic pour l'inclusion : c'est
$Q^1$ qui est une section de $Q^2$}

$\cdot \to H^0(Q^1) \to H^2($ \note{la ligne s'arrête là}

\[
  n = 2\nu \qquad H^2(Q^n) \;\text{---}\; H^{2\nu-2}(Q^n)
\]

\[
  \xi^0 \;-\; \xi^{\nu-1} \quad \xi^{\nu} \quad \xi^{\nu+1} \ \cdots \ \xi^{2\nu}
\]
\note{sous $\xi^{\nu}$, un « $\ell$ » isolé ; au-dessus de $\xi^{\nu+1}$, un
« $[\xi$ » inachevé ; au-dessous, une ligne « $\xi_2^1 \cdots \xi^{\nu}\ \xi$ »
biffée}

\[
  Q^n \quad
  \left\{
  \begin{array}{l}
    \xi^i \in H^{2i}(Q^n, \mu(i)) \qquad \struck{0 \leqslant i \leqslant 2n} \\
    \struck{\lambda_{2\nu} \in H^{2\nu}(Q^{2\nu})}\ \ \lambda_{2\nu} \in H^{2\nu}(Q^{2\nu})
  \end{array}
  \right.
\]

$Q \cdot \xi_i$ \note{isolé, à gauche}

\[
  Q^{n-1} \subset Q^n \qquad
  \xi_i \longleftarrow \xi_i \qquad
  \xi_i \rightsquigarrow \xi_{i+2}
  \qquad
  \begin{pmatrix} 1, & 1 \\ 1, & -1 \end{pmatrix}
\]

\[
  \begin{array}{ccccc}
    Q^0 \subset & Q^1 \subset & Q^2 \subset & Q^3 \subset & Q^4 \\[4pt]
    (\lambda_0, \xi^0) & \xi^0 & \xi^0 & \xi^0 & \xi^0 \\
     & \tfrac12 \xi & (\lambda_1, \xi) & \xi & \xi \\
     & & \tfrac12 \xi^2 & \tfrac12 \xi^2 & (\lambda_2, \xi^2) \\
     & & & \tfrac12 \xi^3 & \tfrac12 \xi^3 \\
     & & & & \tfrac12 \xi^4
  \end{array}
\]
\note{la colonne de $Q^0$ porte « $\lambda_0^{+}\,\xi^0_{+}$ », les indices
douteux}

\struck{$Q^{\ldots}$} : $\xi \ldots$

\[
  \boxed{
  \begin{array}{ll}
    Q^{n = 2\nu} : & \xi^0, \ldots, \xi^{\nu-1}, [\xi^{\nu}, \lambda_{\nu}],
      \tfrac12 \xi^{\nu+1}, \ldots, \tfrac12 \xi^{2\nu} \\[4pt]
    Q^{n = 2\nu+1} : & \xi^0, \ldots, \xi^{\nu}, \tfrac12 \xi^{\nu+1}, \ldots,
      \tfrac12 \xi^{2\nu+1}
  \end{array}}
  \qquad \lambda_{\nu}
\]

\marginal{\ill{} \ill{} $2$ dans $H^{2\nu}(Q^{2\nu})$, \ill{} \ill{}
$\tfrac12(\xi^{\nu} + \lambda_{\nu}) = \varphi_{\nu}$ \ill{} \uncertain{base}
$(\xi^{\nu}, \varphi_{\nu})$}\note{écrit en biais au-dessus du cadre, relié
par un trait à $[\xi^{\nu}, \lambda_{\nu}]$}

\page{24}

\note{feuillet de tapuscrit, paginé \uncertain{« 11 »} : le Lemme 1.14 (si $T$ est
un tore maximal de $G$ et $R$ un sous-tore qui commute à $T$, alors
$R \subset T$), démontré par le sous-groupe image $T'$ de
$R \times_S T \to G$, qui est un tore contenant $T$, donc égal à $T$ ; d'où
1.13 ; le Corollaire 1.15 ($G$ commutatif, lisse et affine : un unique tore
maximal, qui contient tout sous-tore) ; le Corollaire 1.16 ($G$ lisse et affine
à fibres géométriques connexes et nilpotentes : un unique tore maximal, central,
le plus grand sous-tore) et le début de sa preuve, par 1.7 b), 1.13 et
Exp IX 5.6 b), réduite au cas d'un corps algébriquement clos}

Sa main :
\begin{itemize}
  \item les deux premières lignes, tapées et biffées d'un trait ondulé
  (« \struck{En vertu de 10.5.3. c'est un sous-groupe fermé de G lisse} … ») ;
  « Ceci résulte du \struck{résultat suivant} : » ;
  \item la renumérotation à l'encre en 1.4, 1.7, 1.8, 1.13, 1.14, 1.15, 1.16 ;
  \item « Lemme 1.14. Soient $G$ \struck{affine} \add{un $S$-préschéma en
  groupes,} de présentation finie sur $S$ » ;
  \item « En particulier\supplied{, utilisant 6.7. b) :} » ;
  \item Corollaire 1.15 : « commutatif, lisse et affine sur $S$\add{, de rang
  réductif localement constant.} Alors » ;
  \item Corollaire 1.16 : « clôture algébrique \add{$k$} de $k(s)$ » ; la phrase
  « Supposons de plus le rang réductif de $G$ loc. constant. » insérée avant
  « Alors $G$ admet » ; « \add{de plus} $T$ est central » ;
  \item « En vertu de Exp \add{IX}, 5.6. b) », le numéro de l'exposé écrit à
  l'encre sur un autre.
\end{itemize}

\page{25}

\[
  Q^i \subset Q^j \qquad\qquad \xi
\]

\[
  \struck{\mathbf{Z}[T][T'] / T^{2\nu+2},\ 2T' - T^{\nu+1}}
\]

\[
  \left\{
  \begin{array}{l}
    \xi^0, \ldots, \xi^{\nu-1}, [\xi^{\nu}, \lambda_{\nu}], \xi_{\nu+1}, \ldots, \xi_{2\nu} \\
    \xi^0, \ldots, \xi^{\nu}, \xi_{\nu+1}, \ldots, \xi_{2\nu+1}
  \end{array}
  \right.
  \qquad
  \left[
  \begin{array}{l}
    2\xi_i = \xi^i \\
    \xi^i \xi_j = \xi_{i+j} \\
    \xi_i \xi_j = 0 \ \text{pr raisons de degré} \\
    \lambda_{\nu}^2 = 2\xi_{2\nu} = \xi^{2\nu} \\
    \lambda_{\nu} \xi = 0 \\
    \lambda_{\nu} \xi_{\alpha} = 0 \quad \alpha \geqslant \nu+1 \ \text{raisons de degré}
  \end{array}
  \right.
\]
\note{colonne encadrée à droite de la page. « $\lambda_\nu \xi = 0$ » est lu
tel quel}

\[
  Q^i \subset Q^j \quad (i \leqslant j)
\]
\[
  \left.
  \begin{array}{l}
    \xi^{\alpha} \longleftarrow \xi^{\alpha} \\
    \xi_{\alpha} \longleftarrow \xi_{\alpha} \\
    0 \longleftarrow \lambda_{\nu} \\
    \xi_{\nu} \longleftarrow \varphi_{\nu}
  \end{array}
  \right\}\ \text{si}\ i < j = 2\nu
\]
: $\lambda_{\nu}$ \ill{} une classe \uncertain{effective}\note{le premier
« $\longleftarrow$ » est écrit sur une flèche biffée ; l'indice de
$\xi_\nu$ dans la dernière ligne est empâté}

\[
  \left.
  \begin{array}{l}
    \xi^{\alpha} \rightsquigarrow \xi^{\alpha + (j-i)} \\
    \xi_{\alpha} \rightsquigarrow \xi_{\alpha + (j-i)} \\
    \lambda_{\nu} \rightsquigarrow 0 \\
    \varphi_{\nu} \rightsquigarrow \xi_{\nu + (j-i)}
  \end{array}
  \right\}\ \text{si}\ i = 2\nu < j
\]
: $\lambda_{\nu}$ \ill{} une classe « \uncertain{évanescente} ».

\page{26}

\note{feuillet de tapuscrit, paginé « 21 », dans le paragraphe 4 : fin du
critère 4.3 (pour que $Z$ soit un centre réductif de $G$, il faut et il suffit
qu'il soit de type multiplicatif et que chaque $Z_s$ soit un centre réductif de
$G_s$), démontré par Exp IX 5.6 b), 6.8 et 5.1 bis ; la remarque que cette
hypothèse fibre par fibre est purement géométrique ; le Théorème 4.4 (un groupe
algébrique \emph{affine} sur un corps admet un centre réductif) et le début de
sa preuve, par réduction au cas commutatif (Exp VIII 6.7) et algébriquement
clos, où $G = Z \times U$ avec $Z$ de type multiplicatif et $U$ « unipotent »}

Sa main :
\begin{itemize}
  \item « de présentation finie\supplied{, et à fibres affines,} $Z$ un
  sous-préschéma en groupes » ;
  \item la renumérotation à l'encre en 4.2, 4.3, 4.4 ; « Exp IX 5.1. bis » ;
  \item « Nous verrons \struck{\ill{}} \add{dans Exp XIII} que dans ce cas $G$
  s'écrit comme un produit $Z \times U$ » ;
  \item quelques lettres repassées (« En », « de ») et le « $\in$ » ajouté.
\end{itemize}

\page{27}

\note{page écrite en travers de la feuille}

\[
  \begin{bmatrix} 1, & 1 \\ 1, & -1 \end{bmatrix}
\]

\[
  \begin{array}{ccccccc}
    Q^0 & Q^1 & Q^2 & Q^3 & Q^4 & Q^5 & Q^6 \\[4pt]
    [\lambda_0, \xi^0] & \xi^0 = 1 & \xi^0 = 1 & \xi^0 = 1 & \xi^0 = 1 & \xi^0 = 1 & \xi^0 = 1 \\
     & \tfrac12 \xi & [\lambda_1, \xi] & \xi & \xi & \xi & \xi \\
     & & \tfrac12 \xi^2 & \tfrac12 \xi^2 & [\lambda_2, \xi^2] & \xi^2 & \xi^2 \\
     & & & \tfrac12 \xi^3 & \tfrac12 \xi^3 & \tfrac12 \xi^3 & [\lambda_3, \xi^3] \\
     & & & & \tfrac12 \xi^4 & \tfrac12 \xi^4 & \tfrac12 \xi^4 \\
     & & & & & \tfrac12 \xi^5 & \tfrac12 \xi^5 \\
     & & & & & & \tfrac12 \xi^6
  \end{array}
\]
\note{colonnes réunies ici en un tableau ; sur la page chaque quadrique a sa
colonne. Dans celle de $Q^1$, le $\xi$ porte une marque empâtée en exposant}

\struck{a) Cohomologie de $Q^n$ nulle en dim impaire}\note{au pied de la page,
tête-bêche et biffé : c'est le a) de la page 21}

\page{28}

\note{feuillet de tapuscrit, paginé « 12 » : fin de la preuve de 1.16 ($T(k)$
est dans le centre de $G(k)$ par BIBLE 6 th. 2, donc $T$ est dans le centre de
$G$) ; la Proposition 1.17, « triviale », sur $G \supset H \supset T$ ; puis le
début du paragraphe 2, « Le groupe de Weyl » : sur un corps algébriquement
clos, $C$ et $N$ centralisateur et normalisateur d'un tore maximal $T$ sont des
sous-groupes fermés lisses, $C$ ouvert dans $N$, $W = N/C$ est un groupe fini
étale, le « groupe de Weyl géométrique » de $G$ relativement à $T$ ; par le
théorème de conjugaison de Borel il ne dépend pas du tore maximal, et sa
formation commute à l'extension de la base}

Sa main :
\begin{itemize}
  \item « ce qui implique que $T$ est dans le centre de $G$, \struck{en vertu de
  Exp VI} \add{car $\mathrm{Centr}_G(T)$ est un sous-schéma fermé de $G$
  (Exp VIII 6.6.), qui contient les points de $G(k)$, donc est identique à $G$
  par le Nullstellensatz, ($G$ étant réduit).} »
  \marginal{Nullstellensatz}\note{entouré, dans la marge gauche} ;
  \item la renumérotation en 1.17 et 2. ;
  \item Proposition 1.17 : « Soient $G \supset H \supset T$ des préschémas en
  groupes de type fini sur $S$. \struck{Alors} \add{Conditions équivalentes :}
  (i) $T$ est un tore maximal de $G$\struck{, c'est un tore maximal de $H$, et
  pour tout $s \in S$, $G_s$ et $H_s$ ont même rang réductif} (ii) $T$ est un
  tore maximal de $H$, et pour tout $s \in S$, on a égalité des rangs réductifs
  $\rho_r(G_s) = \rho_r(H_s)$\struck{, alors $T$ est un tore maximal de $G$}. »,
  les « (i) », « (ii) » et les $\rho_r$ étant de sa main ;
  \item « soit $G$ un groupe algébrique sur $k$\add{, lisse et affine sur
  $k$.} Si $T$ est un tore maximal » ; « en vertu de \add{Exp XI} 5.9. » ;
  « \add{$W =$} $N/C$ » ; deux passages tapés biffés sur la définition du
  groupe de Weyl ; « (ou simplement groupe de Weyl si une confusion n'est pas à
  craindre), de $G$ relativement à $T$ » repris entre parenthèses et
  soulignés ; « du théorème de \add{conjugaison de} Borel ».
\end{itemize}

\page{29}

\note{les deux lignes qui suivent sont au crayon}

\struck{$\boxed{F \otimes G}$}

\[
  L_{\cdot}(F_{\cdot}) \otimes L_{\cdot}(G_{\cdot}) = F \mathbin{\underline{\otimes}} G
\]
\note{le second $\otimes$ est souligné deux fois}

\[
  n = 2\nu
\]

\[
  H^{n-1}(U^{n-1})(1) \xrightarrow[\partial]{\text{ins}} H^n(Q^n)
  \xrightarrow{\text{res}} H^n(U^n)
\]
\note{l'étiquette « ins » de la première flèche est une lecture douteuse. Une flèche courbe va du premier terme au troisième, marquée
« \struck{\ill{}} bij » et d'une fraction
« $\mathrm{exc}^2 / \ell^{(2\nu+1)} \rightsquigarrow 2\ell^{(2\nu)}$ », lecture
douteuse}

\[
  \begin{array}{lll}
    \ell^{2\nu+1} & & \\
    \ell^{2\nu+1} \rightsquigarrow \lambda_{\nu} &
    [\lambda_{\nu}, \xi^{\nu}] \rightsquigarrow [2\ell^{2\nu}, 0] & \\
     & \varphi_{\nu} = \tfrac12 (\lambda_{\nu} + \xi^{\nu}) \rightsquigarrow \ell^{2\nu} &
  \end{array}
\]
\note{le premier « $\ell^{2\nu+1}$ » est écrit sur un autre exposant}

\struck{$H^0(Q^0) \to H^0(U^0)$}

\[
  \lambda_{\nu} \longleftarrow \ell_{2\nu}
\]

\[
  H^{n+1}_c(U^{n+1})(+1) \xleftarrow[\partial]{} H^n(Q^n)
  \xleftarrow{i_{\partial}} H^n_c(U^n)
\]

\note{au-dessus de la première flèche, une étiquette illisible}

\[
  \begin{array}{ll}
    [2\ell_{2\nu+1}, 0] \longleftarrow [\lambda_{\nu}, \xi^{\nu}] & \\
    \ell_{2\nu+1} \longleftarrow \varphi_{\nu} & \\
    2\ell_{2\nu+1} \longleftarrow \ell_{2\nu} &
  \end{array}
\]
\note{la dernière ligne est portée sur une flèche courbe du dernier terme de la
suite au premier ; une flèche verticale, à droite, relie la suite du haut à
celle du bas}

\[
  Q^{n+1} - Q^{n} = U^{n+1}
\]

\emph{N.B.} On \uncertain{définit} $\varphi_{\nu}$
\emph{\uncertain{directement}} par un \ill{} \uncertain{isotrope}
\uncertain{maximal} — les deux familles de \uncertain{tels} \ill{} \ill{} des
classes $\alpha_{\nu}$, $\beta_{\nu}$. \ill{} \uncertain{composantes} : on a
\[
  \alpha_{\nu} + \beta_{\nu} = \xi^{\nu}, \qquad
  \alpha_{\nu} - \beta_{\nu} = \lambda_{\nu}, \qquad
  \varphi_{\nu} = \alpha_{\nu} \ \ldots
\]

\page{30}

\note{feuillet de tapuscrit, paginé « 26 » : démonstration de a) et b) d'un
énoncé sur le centre réductif — pour a) et la première assertion de b), renvoi à
1.7 et à Exp IX 5.6 ; puis la seconde assertion de b), que $G' = G/Z$ a pour
centre réductif le sous-groupe unité : $G/Z$ est représentable, affine, de
présentation finie et lisse sur $S$ (Exp VIII 5.1, 5.8, Exp VI), à fibres
connexes ; on se ramène au cas d'un corps (4.3) et on fait opérer l'image
inverse $Z_1$ d'un sous-groupe central $Z'$ de type multiplicatif de $G'$ par
automorphismes intérieurs. De sa main, seulement un crochet devant « sont
également triviales »}

\page{31}

\[
  Q^{2\nu} : \qquad \xi : H^{2i}(Q^{2\nu}) \to H^{2i+2}(Q^{2\nu})(1)
  \qquad \struck{\ill{}}\ (0 \leqslant i \leqslant 2\nu - 1)
\]
\note{devant chacun des deux $Q^{2\nu}$, un signe biffé ; la borne de
droite est surchargée}

\struck{injectif si $i < 2\nu$ i.e. $i \leqslant \nu - 1$}

\emph{bijectif} si $i \neq \nu - 1, \nu$

$i = \nu - 1$ : \emph{injectif} \struck{si $i = \nu - 1$}, conoyau libre de base
$\varphi_{\nu}$

$i = \nu$ : \struck{\uncertain{noyau}} \emph{surjectif}, noyau libre de base
$\lambda_{\nu}$

\[
  Q^{2\nu+1} \qquad \xi : H^{2i}(Q^{2\nu+1}) \to H^{2i+2}(Q^{2\nu+1})(1)
  \qquad (0 \leqslant i \leqslant 2\nu)
\]
\note{le « $+1$ » du dernier exposant est entouré, ajouté au-dessus}

\emph{bijectif} si $i \neq \nu$

injectif de conoyau $\mathbf{Z}/2\mathbf{Z}$ si $i = \nu$

\page{32}

\note{feuillet de tapuscrit, paginé « 31 » : suite de la démonstration d'un
énoncé sur le rang réductif $\rho_r$ — fin de a) ($\dim C' \leqslant \dim C$) ;
b) si $\rho_r$ est localement constant, un tore maximal dans une fibre l'est au
voisinage ; les équivalences (i) $\Leftrightarrow$ (ii bis) et la réciproque
(i) $\Rightarrow$ (ii) par le foncteur $\mathcal{T}$ des tores maximaux, sous-foncteur
du foncteur $F$ de Exp XI 4.1, représentable par un sous-préschéma ouvert de
$F$, donc lisse, et admettant une section localement pour la topologie étale
(Exp XI 1.10) ; c) conséquence de Exp XI 5.4 bis et du théorème de conjugaison
de Borel ; début de d)}

Sa main :
\begin{itemize}
  \item « donc $\dim C' \leqslant \dim C$ », le signe à l'encre ; une ligne tapée
  biffée au début de b) ;
  \item « si $T_s$ est \add{tore} maximal dans $G_s$ » ; les $\rho_r$ et les
  flèches $\Rightarrow$, $\Leftrightarrow$ à l'encre ;
  \item « reste à prouver l'implication inverse \add{(i) $\Rightarrow$ (ii)} » ;
  \item « le foncteur $F$ de \add{Exp XI} 4.1., qui est donc un préschéma
  \struck{\ill{}} lisse \add{et séparé} sur $S$ » ;
  \item le sous-foncteur noté $\mathcal{T}$ à l'encre, et
  \marginal{$\mathcal{T}$}\note{dans la marge gauche} ;
  « pour un $S'$ sur $S$ », « $G_{S'}$ », « point $s \in S'$ » corrigés à
  l'encre ; « Utilisant la remarque\supplied{, moyennant (i),} que » ;
  \item « et par suite est lisse \add{et séparé} sur $S$ » ; « une section
  \struck{localement} localement au sens de la topologie étale en vertu de
  \add{Exp XI} 1.10. » ; « c) C'est une conséquence immédiate de \add{Exp XI}
  5.4. bis ».
\end{itemize}

\page{34}

\emph{Quadriques} Une base pour l'équivalence linéaire dans $Q^{2n}$ est formée
de
\[
  1, H, H^2, \ldots, H^{n-1}, \underbrace{S', S''}, HS' = HS'',
  \ldots, H^n S' = H^n S''
\]
\note{sous l'accolade, un mot illisible}

On a $H^n = S' + S''$

[$H$ est une section hyperplane dans l'immersion projective évidente, $S'$ et
$S''$ sont définis par des sous-espaces \struck{linéaires de dimension}
\uncertain{de dim.} \ill{} $Q^{2n}$, appartenant aux deux familles algébriques
distinctes de tels sous-espaces]

Une base de $Q^{2n+1}$ est formée de
\[
  1, H, \ldots, H^n, \frac{H^{n+1}}{2}, \ldots, \frac{H^{2n+1}}{2}
\]
[où $H$ est une section hyperplane, et où les formes
$\frac{H^{n+1}}{2} \ldots \frac{H^{2n+1}}{2}$ correspondent aux sous-variétés
linéaires contenues dans $Q^{2n+1}$]

\note{feuillet de papier bistre, d'une écriture plus posée que celle des pages
qui précèdent ; les mots courts restent réduits à des traits et sont lus
d'après le sens}

\section*{Le tiré à part}

\page{35}

\note{première page d'un tiré à part imprimé : Antoine Delzant, « Définition
des classes de Stiefel-Whitney d'un module quadratique sur un corps de
caractéristique différente de 2 », Note transmise par M. Jean Leray, extraite
des \emph{Comptes rendus des séances de l'Académie des Sciences}, t. 255,
p. 1366-1368, séance du 10 septembre 1962 (d'après la page 38, qui ne porte
que cette mention). Pour $A$ de caractéristique $\neq 2$, l'anneau $Q(A)$ des
classes de formes quadratiques, engendré par le monoïde des classes de modules
quadratiques non dégénérés, et son augmentation (le rang) ; $D(A) = A^*/2A^*$
identifié à $H^1(A, \mathbf{Z}/2\mathbf{Z})$ ; l'homomorphisme
$\varphi : \mathbf{Z}(D) \to Q(A)$ ; la Proposition : $\varphi$ est surjectif,
compatible aux augmentations, de noyau engendré par les
$e_a + e_b - e_{a'} - e_{b'}$ tels que $\varphi_a + \varphi_b = \varphi_{a'} +
\varphi_{b'}$ ; le cup-produit $D(A) \times D(A) \to H^2(A, \mathbf{Z}/2\mathbf{Z})$,
$(e_a, e_b) \mapsto a.b$, vu dans le groupe de Brauer. Rien de sa main}

\page{36}

\note{page 2 du tiré à part : $a.b$ est la classe de l'algèbre de quaternions
$(a, b)$ ; par un théorème de Witt, le noyau de $\varphi$ est engendré par les
$e_a + e_b - e_{a'} - e_{b'}$ avec $e_{a+b} = e_{a'+b'}$ et $a.b = a'.b'$ ; pour
$H$ un anneau gradué et $\widetilde{H} = \mathbf{Z} \times (1 + \hat H^+)$
« l'anneau filtré déduit de $H$ par le procédé exposé par Grothendieck » (Bull.
Soc. math. France 86, 1958), un homomorphisme $w^1 : D(A) \to H^1$ vérifiant la
condition du corollaire définit $Q(A) \to \widetilde{H}$ ; avec $H = H(A)$ la
cohomologie galoisienne à coefficients $\mathbf{Z}/2\mathbf{Z}$, on obtient
$w : Q(A) \to \widetilde{H(A)}$ et les classes de Stiefel-Whitney
$w^i(q) \in H^i(A, \mathbf{Z}/2\mathbf{Z})$, avec
$w^n(q + q') = \sum_{i+j=n} w^i(q) w^j(q')$ ; pour $q = \sum a_j x_j^2$,
$w^1(q) = \sum a_j$ (le discriminant) et $w^2(q) = \sum_{i<j} a_i . a_j$,
interprété par les algèbres de Clifford (Springer) ; $w$ est un homomorphisme
de foncteurs. Rien de sa main}

\page{37}

\note{page 3 du tiré à part : sur un corps de nombres, l'application
$Q(A) \to \prod_v Q(A_v)$ est injective (Hasse) ; sur un complété discret une
forme est déterminée par $\varepsilon$, $w^1$, $w^2$ ; sur $\mathbf{R}$,
$H(\mathbf{R}) = \mathbf{Z}/2\mathbf{Z}[t]$ et $w^j(q) = \binom{n}{j} \bmod 2$,
$n$ le nombre de carrés négatifs ; d'où, par un diagramme commutatif, que sur
un corps de nombres une forme quadratique est déterminée par son rang et ses
classes de Stiefel-Whitney. Remarque : on ne sait pas si
$\mathrm{grad}^2(Q(A)) \to H^2(A, \mathbf{Z}/2\mathbf{Z})$ est toujours injective,
ni si tout élément d'ordre 2 du groupe de Brauer est somme d'éléments $a.b$.
Références : Bourbaki, \emph{Algèbre} IX ; Grothendieck, Bull. SMF 86 ;
Springer, Proc. Akad. Amst. 62 A ; Witt, Crelle 176. Rien de sa main}

\end{document}
